\chapterauthor{Lech Tuzinkiewicz}
\chapter[Projektowanie baz \ldots]{Projektowanie baz danych sterowane jakością}
\thispagestyle{empty}
\vspace{-8mm} {\large Lech Tuzinkiewicz} \vspace{8mm}

\section{Wprowadzenie}
Obecnie działalność naukowa i gospodarka, zdominowane są przez procesy cyfryzacji i transformacji. Na naszych oczach dokonuje się fundamentalna zmiana paradygmatu danych, przy czym paradygmat należy rozumieć jako \textbf{koncepcję, perspektywę lub zbiór idei, które mają wpływ na sposób rozumienia i interpretacji wycinka rzeczywistości w określonej dziedzinie}. W tym kontekście dane, które jeszcze kilkanaście lat temu traktowano jako element działalności operacyjnej, obecnie uważane są za produkt o znacznej wartości i mający ogromne znaczenie w kształtowaniu strategii konkurencyjnej. \textbf{Dane} nie są już jedynie zasobem gromadzonym w bazach danych i wykorzystywanym w operacjach biznesowych, ale stanowią zarówno źródło pozyskiwania informacji, jak i są istotnym zasobem majątkowym czyli walorem w postaci \textbf{aktywu }o znaczeniu często przewyższającym tradycyjne aktywa materialne, takie jak środki trwałe czy kapitał finansowy. Aktualnie również uważa się, że dane można postrzegać jako \textbf{najcenniejszy aktyw we współczesnym świecie.}

W erze sztucznej inteligencji i uczenia maszynowego, jakościowo dobre dane mają istotne znaczenie dla osób i systemów informatycznych wykorzystywanych w podejmowaniu decyzji strategicznych w oparciu o analitykę. Znaczenie \textbf{jakości danych} w kontekście ich wartości jest absolutnie \textbf{fundamentalne. }Wyraża się to, między innymi powszechnym stwierdzeniem w postaci GIGO tzn. „Garbage In, Garbage Out”, czyli jakość podejmowanych decyzji odpowiada jakości danych, które w danym procesie są wykorzystywane. Podsumowując, jakość danych ma wpływ na decyzje biznesowe oraz modele AI i tym samym mogą być one tylko tak dobre, jak dane, na których się opierają.
Dane gromadzone są w różnych bazach danych (np. relacyjnych, NoSQL, Dokumentowych itd.) i proces ich projektowania musi uwzględniać aspekt oceny jakości danych, które mają być gromadzone i udostępniane w trakcie realizacji procesów biznesowych. 

W dalszej części rozdziału zostanie przedstawiony krótki przegląd literatury obejmujący wybrane aspekty oceny jakościowej procesu projektowania baz danych oraz przykładowe modele jakości mające zastosowanie w tym procesie. 


